\section{Conclusion}
This paper proposes AgentCAT, a Large Language Model-based multi-agent collaborative simulation system designed to break through the constraints of traditional CAT research, which is chronically limited by static offline data and isolated optimization. Extensive experimental analysis robustly demonstrates the effectiveness and robustness of this simulation system in both ability estimation and selection strategies.
 
However, despite AgentCAT's superior overall performance, its current architecture still exhibits certain limitations, specifically in the design of the supervisor. Currently, when executing ability updates, the supervisor primarily relies on pre-defined mathematical rules. While this formulaic update paradigm is statistically interpretable, it may cause ability estimation to fall into local optima when faced with accumulated noise from long-term interactions or atypical response behaviors of the examinee. Future work will be dedicated to introducing Reinforcement Learning mechanisms to reconstruct the ability estimation module. We plan to design an IRT-based ability estimation agent to replace traditional formulaic numerical solvers. This agent will not rely on static mathematical formulas, but will instead model the ability update process as a sequential decision-making problem. It will be capable of intelligently perceiving subtle fluctuations in the examinee's cognitive state based on the current dialogue history and exercise features, dynamically adjusting the step size and direction of ability updates. Through this approach, we expect to achieve a paradigm shift in ability assessment from mechanical calculation to intelligent perception, thereby further enhancing the system's adaptability and assessment precision in complex, uncertain educational environments.